\documentclass[12pt,a4paper]{article}
\usepackage[indonesian]{babel}
\usepackage[utf8]{inputenc}
\usepackage{geometry}
\usepackage{graphicx}
\usepackage{longtable}
\usepackage{booktabs}
\usepackage{array}
\usepackage{xcolor}
\usepackage{listings}
\usepackage{hyperref}
\usepackage{fancyhdr}
\usepackage{titlesec}
\usepackage{tocloft}

% Page setup
\geometry{margin=2.5cm}
\pagestyle{fancy}
\fancyhf{}
\fancyhead[L]{Laporan Analisis Sistem}
\fancyhead[R]{Samak Pro}
\fancyfoot[C]{\thepage}

% Colors
\definecolor{codegreen}{rgb}{0,0.6,0}
\definecolor{codegray}{rgb}{0.5,0.5,0.5}
\definecolor{codepurple}{rgb}{0.58,0,0.82}
\definecolor{backcolour}{rgb}{0.95,0.95,0.92}

% Code listing style
\lstdefinestyle{mystyle}{
    backgroundcolor=\color{backcolour},   
    commentstyle=\color{codegreen},
    keywordstyle=\color{blue},
    numberstyle=\tiny\color{codegray},
    stringstyle=\color{codepurple},
    basicstyle=\ttfamily\footnotesize,
    breakatwhitespace=false,         
    breaklines=true,                 
    captionpos=b,                    
    keepspaces=true,                 
    numbers=left,                    
    numbersep=5pt,                  
    showspaces=false,                
    showstringspaces=false,
    showtabs=false,                  
    tabsize=2
}
\lstset{style=mystyle}

% Title formatting
\titleformat{\section}{\Large\bfseries}{BAB \thesection:}{0.5em}{}
\titleformat{\subsection}{\large\bfseries}{\thesubsection}{0.5em}{}

\begin{document}

% ============ COVER PAGE ============
\begin{titlepage}
    \centering
    \vspace*{2cm}
    
    {\Huge\bfseries LAPORAN ANALISIS SISTEM INFORMASI\par}
    \vspace{0.5cm}
    {\Huge\bfseries MANAJEMEN MASJID\par}
    \vspace{1.5cm}
    
    {\LARGE\bfseries Samak Pro\par}
    {\Large Aplikasi Manajemen Masjid\par}
    
    \vspace{2cm}
    
    \includegraphics[width=0.3\textwidth]{example-image}
    
    \vspace{2cm}
    
    {\Large Disusun oleh:\par}
    {\large Tim Pengembang Samak Pro\par}
    
    \vfill
    
    {\large Politeknik Negeri Bandung\par}
    {\large Proyek 3\par}
    {\large Desember 2025\par}
\end{titlepage}

% ============ TABLE OF CONTENTS ============
\tableofcontents
\newpage

% ============ BAB 1 ============
\section{RANCANGAN}

\subsection{Latar Belakang}
Masjid sebagai pusat kegiatan umat Islam memerlukan pengelolaan yang efektif dan efisien. Pengelolaan manual sering mengalami kendala seperti pencatatan keuangan yang tidak transparan, informasi kegiatan yang tidak tersebar luas, dan kesulitan dalam berinteraksi dengan jamaah. \textbf{Samak Pro} hadir sebagai solusi digital untuk membantu pengurus DKM (Dewan Kemakmuran Masjid) mengelola berbagai aspek kegiatan masjid secara terintegrasi.

\subsection{Tujuan Aplikasi}
\begin{enumerate}
    \item \textbf{Digitalisasi Pengelolaan Masjid} -- Mengubah proses manual menjadi digital terkomputerisasi
    \item \textbf{Transparansi Keuangan} -- Menyediakan laporan keuangan yang dapat diakses publik
    \item \textbf{Penyebaran Informasi} -- Mempermudah publikasi berita, jadwal kajian, dan kegiatan masjid
    \item \textbf{Layanan Jamaah} -- Menyediakan fitur konsultasi online dan layanan barang hilang
    \item \textbf{Pengelolaan Donasi} -- Memfasilitasi donasi online dengan perhitungan zakat otomatis
\end{enumerate}

\subsection{Ruang Lingkup}
\begin{table}[h]
\centering
\begin{tabular}{|l|l|}
\hline
\textbf{Aspek} & \textbf{Cakupan} \\
\hline
Domain & Manajemen Masjid \& Layanan Jamaah \\
Platform & Aplikasi Web (Browser-based) \\
Pengguna & Pengurus DKM \& Jamaah Umum \\
Fitur Utama & Postingan, Keuangan, Donasi, Konsultasi, Lost \& Found \\
\hline
\end{tabular}
\caption{Ruang Lingkup Aplikasi}
\end{table}

\subsection{Definisi Pengguna (User Roles)}
\begin{longtable}{|l|l|p{6cm}|c|}
\hline
\textbf{Role} & \textbf{Alias} & \textbf{Deskripsi} & \textbf{Dashboard} \\
\hline
\endfirsthead
\hline
\textbf{Role} & \textbf{Alias} & \textbf{Deskripsi} & \textbf{Dashboard} \\
\hline
\endhead
Super Admin & Super Administrator & Akses penuh untuk development & Ya \\
Admin & Admin System & Mengelola pengguna \& roles & Ya \\
Koordinator Humas & Koor Humas & Menyetujui/menolak postingan & Ya \\
Humas & Bidang Humas & Mengelola konten \& konsultasi & Ya \\
Bendahara Pemasukan & Bendahara M & Mengelola keuangan masuk & Ya \\
Bendahara Pengeluaran & Bendahara K & Mengelola keuangan keluar & Ya \\
Sarpras & Divisi Sarpras & Mengelola barang hilang & Ya \\
Jamaah & Jamaah & Pengguna umum & Tidak \\
\hline
\caption{Daftar Role Pengguna}
\end{longtable}

\newpage

% ============ BAB 2 ============
\section{REQUIREMENT ANALISIS}

\subsection{Kebutuhan Fungsional}

\subsubsection{Modul Autentikasi (AuthController)}
\begin{longtable}{|l|l|l|p{4cm}|}
\hline
\textbf{Fitur} & \textbf{Method} & \textbf{Route} & \textbf{Deskripsi} \\
\hline
\endfirsthead
\hline
\textbf{Fitur} & \textbf{Method} & \textbf{Route} & \textbf{Deskripsi} \\
\hline
\endhead
Login & login() & /login & Masuk ke sistem \\
Register & register() & /register & Pendaftaran baru \\
Logout & logout() & /logout & Keluar dari sistem \\
Kirim OTP & sendOtp() & /auth/send-otp & Verifikasi email \\
Verifikasi OTP & verifyOtp() & /auth/verify & Konfirmasi kode OTP \\
Forgot Password & sendResetLink() & /forgot-password & Reset password \\
\hline
\caption{Fitur Modul Autentikasi}
\end{longtable}

\subsubsection{Modul Manajemen Pengguna (UserController)}
\begin{longtable}{|l|l|l|l|}
\hline
\textbf{Fitur} & \textbf{Method} & \textbf{Route} & \textbf{Permission} \\
\hline
Daftar Pengguna & index() & /admin/users & view\_users \\
Tambah Pengguna & store() & /admin/users/create & edit\_users \\
Edit Pengguna & update() & /admin/users/\{id\}/edit & edit\_users \\
Hapus Pengguna & destroy() & /admin/users/\{id\} & delete\_users \\
\hline
\caption{Fitur Modul Manajemen Pengguna}
\end{longtable}

\subsubsection{Modul Postingan (PostinganController)}
\begin{longtable}{|l|l|l|}
\hline
\textbf{Fitur} & \textbf{Route} & \textbf{Permission} \\
\hline
Daftar Postingan (Client) & /postingan & Public \\
Detail Postingan & /postingan/\{slug\} & Public \\
Daftar Admin & /admin/postingan & view\_posts \\
Tambah Postingan & /admin/postingan/tambah & create\_posts \\
Edit Postingan & /admin/postingan/edit/\{id\} & edit\_posts \\
Hapus Postingan & /admin/postingan/delete/\{id\} & delete\_posts \\
Approval Postingan & /admin/postingan/approval & approve\_posts \\
\hline
\caption{Fitur Modul Postingan}
\end{longtable}

\subsubsection{Modul Galeri (GaleriController)}
\begin{longtable}{|l|l|l|}
\hline
\textbf{Fitur} & \textbf{Route} & \textbf{Permission} \\
\hline
Galeri Publik & /galeri & Public \\
Daftar Admin & /admin/galeri & view\_gallery \\
Tambah Album & /admin/galeri/tambah & create\_gallery \\
Edit Album & /admin/galeri/edit/\{id\} & edit\_gallery \\
Hapus Album & /admin/galeri/delete/\{id\} & delete\_gallery \\
\hline
\caption{Fitur Modul Galeri}
\end{longtable}

\subsubsection{Modul Konsultasi}
\textbf{Client (ConsultationClientController):}
\begin{itemize}
    \item Landing Konsultasi -- /konsultasi
    \item Buat Konsultasi -- /konsultasi-saya/buat
    \item Riwayat -- /konsultasi-saya/riwayat
    \item Detail Chat -- /konsultasi-saya/\{id\}
    \item Kirim Pesan -- /konsultasi-saya/\{id\}/pesan
\end{itemize}

\textbf{Admin/Ustadz (ConsultationUstadzController):}
\begin{itemize}
    \item Daftar Konsultasi (permission: view\_consultations)
    \item Terima/Tolak Konsultasi (permission: reply\_consultations)
    \item Tutup Konsultasi (permission: reply\_consultations)
\end{itemize}

\subsubsection{Modul Donasi/ZIS (ZISController)}
\begin{longtable}{|l|l|p{5cm}|}
\hline
\textbf{Fitur} & \textbf{Route} & \textbf{Deskripsi} \\
\hline
Informasi Donasi & /donasi & Halaman informasi donasi \\
Form Donasi & /donasi/sekarang & Form perhitungan zakat/infaq \\
Hitung Donasi & /donasi/hitung & Proses kalkulasi \\
Konfirmasi & /donasi/konfirmasi & Konfirmasi pembayaran \\
Sukses & /donasi/sukses & Halaman sukses \\
\hline
\caption{Fitur Modul Donasi}
\end{longtable}

\subsubsection{Modul Keuangan (KeuanganController)}
\begin{longtable}{|l|l|l|}
\hline
\textbf{Fitur} & \textbf{Method} & \textbf{Permission} \\
\hline
Dashboard Keuangan & index() & view\_finance \\
Tambah Transaksi & store() & manage\_income / manage\_expense \\
Hapus Transaksi & destroy() & manage\_income \\
Laporan Publik & clientIndex() & Public \\
\hline
\caption{Fitur Modul Keuangan}
\end{longtable}

\subsubsection{Modul Barang Hilang \& Ditemukan}
\begin{itemize}
    \item \textbf{FoundItem (LostFoundController):} CRUD barang ditemukan
    \item \textbf{LostItem (LostItemController):} CRUD laporan barang hilang
    \item Permission: view\_lost\_items, create\_lost\_items, edit\_lost\_items, delete\_lost\_items
\end{itemize}

\subsubsection{Modul Jadwal Kegiatan}
\begin{itemize}
    \item \textbf{JadwalKegiatanController:} Tampilan kalender publik, detail kegiatan
    \item \textbf{AdminKegiatanController:} CRUD kegiatan (permission: view\_events, create\_events, edit\_events, delete\_events)
\end{itemize}

\subsubsection{Modul Form Builder (FormBuilderController)}
\begin{itemize}
    \item Buat Form -- Membuat form dinamis dengan berbagai tipe field
    \item Kelola Response -- Melihat dan menghapus respons form
    \item Publik -- Pengisian form oleh publik
\end{itemize}

\subsubsection{Modul Bank \& Infaq}
\begin{itemize}
    \item \textbf{BankController:} CRUD rekening bank untuk donasi
    \item \textbf{InfaqController:} CRUD program infaq
    \item \textbf{DonationConfirmationController:} Verifikasi donasi online
    \item \textbf{OfflineDonationController:} Input donasi offline
    \item \textbf{KotakAmalController:} Pendataan kotak amal
\end{itemize}

\subsection{Kebutuhan Non-Fungsional}

\subsubsection{Halaman Tampilan}
\begin{table}[h]
\centering
\begin{tabular}{|l|p{10cm}|}
\hline
\textbf{Kategori} & \textbf{Halaman} \\
\hline
Publik & Home, Berita, Galeri, Jadwal, Donasi, Laporan Keuangan, Konsultasi, Form, Tentang Kami \\
\hline
Autentikasi & Login, Register, Verify OTP, Forgot Password, Reset Password \\
\hline
Admin & Dashboard, User/Role, Postingan, Galeri, Kegiatan, Keuangan, Donasi, Konsultasi, Lost \& Found \\
\hline
Profil & Edit Profil, Ubah Password \\
\hline
\end{tabular}
\caption{Daftar Halaman Tampilan}
\end{table}

\subsubsection{Keamanan \& Autentikasi}
\begin{table}[h]
\centering
\begin{tabular}{|l|p{8cm}|}
\hline
\textbf{Fitur} & \textbf{Implementasi} \\
\hline
Autentikasi & Laravel Auth dengan Email Verification \\
Otorisasi & Role-Based Access Control (RBAC) dengan Permission \\
Middleware & CheckPermission, CheckRole, auth \\
Verifikasi Email & OTP via Email \\
Password Reset & Token-based via Email \\
Captcha & Trait VerifyCaptcha untuk registrasi \\
\hline
\end{tabular}
\caption{Fitur Keamanan}
\end{table}

\newpage

% ============ BAB 3 ============
\section{SYSTEM DESIGN (Perancangan Sistem)}

\subsection{Stack Teknologi}
\begin{table}[h]
\centering
\begin{tabular}{|l|l|l|}
\hline
\textbf{Layer} & \textbf{Teknologi} & \textbf{Versi} \\
\hline
Backend Framework & Laravel & 12.x \\
Bahasa Pemrograman & PHP & >= 8.2 \\
Database & PostgreSQL & >= 14 \\
Frontend & Blade Template + Tailwind CSS & 4.x \\
Build Tool & Vite & 7.x \\
Real-time & Laravel Reverb + Pusher & - \\
Package Manager & Composer + NPM & - \\
\hline
\end{tabular}
\caption{Stack Teknologi}
\end{table}

\subsection{Alur Teknologi (MVC Pattern)}
Laravel menggunakan arsitektur MVC (Model-View-Controller):

\begin{enumerate}
    \item \textbf{Request} -- Browser mengirim HTTP request ke server
    \item \textbf{Routing} -- web.php mencocokkan URL dengan controller
    \item \textbf{Middleware} -- CheckPermission \& auth memvalidasi akses
    \item \textbf{Controller} -- Memproses logika bisnis
    \item \textbf{Model} -- Berinteraksi dengan database via Eloquent ORM
    \item \textbf{View} -- Blade template merender HTML
    \item \textbf{Response} -- HTML dikirim ke browser
\end{enumerate}

\subsection{Database Schema}
\begin{longtable}{|l|p{8cm}|}
\hline
\textbf{Tabel} & \textbf{Deskripsi} \\
\hline
\endfirsthead
users & Data pengguna \\
roles & Daftar role \\
permissions & Daftar permission \\
role\_permission & Pivot role-permission \\
postingans & Artikel/berita \\
gallery\_albums & Album galeri \\
gallery\_photos & Foto dalam album \\
jadwal\_kegiatans & Jadwal kegiatan \\
consultations & Konsultasi jamaah \\
consultation\_messages & Pesan chat konsultasi \\
bank\_accounts & Rekening bank \\
donation\_confirmations & Konfirmasi donasi \\
financial\_transactions & Transaksi keuangan \\
found\_items & Barang ditemukan \\
lost\_items & Laporan barang hilang \\
forms & Form builder \\
form\_fields & Field dalam form \\
form\_responses & Respons form \\
infaqs & Program infaq \\
kotak\_amals & Pendataan kotak amal \\
website\_informations & Halaman statis \\
\hline
\caption{Database Schema}
\end{longtable}

\subsection{Arsitektur Aplikasi}
Aplikasi Samak Pro menggunakan arsitektur monolitik dengan Laravel sebagai backend framework. Komponen utama:

\begin{itemize}
    \item \textbf{Frontend:} Blade Views, Tailwind CSS, JavaScript/Alpine
    \item \textbf{Backend:} Laravel 12, Eloquent ORM, Queue System
    \item \textbf{Services:} Email Service, File Storage, WebSocket (Reverb)
    \item \textbf{Database:} PostgreSQL
\end{itemize}

\subsection{Interfaces}

\subsubsection{Hardware Interface}
\begin{itemize}
    \item Server dengan minimal 2GB RAM
    \item Storage minimal 20GB untuk file upload
    \item Koneksi internet stabil
\end{itemize}

\subsubsection{Software Interface}
\begin{table}[h]
\centering
\begin{tabular}{|l|l|}
\hline
\textbf{Komponen} & \textbf{Interface} \\
\hline
Web Server & Nginx / Apache \\
PHP Runtime & PHP-FPM >= 8.2 \\
Database & PostgreSQL Driver (pdo\_pgsql) \\
Cache & File / Redis \\
Queue & Database / Redis \\
\hline
\end{tabular}
\caption{Software Interface}
\end{table}

\subsubsection{Komunikasi (Protokol)}
\begin{table}[h]
\centering
\begin{tabular}{|l|p{8cm}|}
\hline
\textbf{Protokol} & \textbf{Penggunaan} \\
\hline
HTTP/HTTPS & Request-response web \\
WebSocket & Real-time chat konsultasi (Laravel Reverb) \\
SMTP & Pengiriman email (OTP, reset password) \\
REST & API internal untuk AJAX \\
\hline
\end{tabular}
\caption{Protokol Komunikasi}
\end{table}

\newpage

% ============ BAB 4 ============
\section{UI (User Interface)}

\subsection{Struktur Halaman}

\subsubsection{Halaman Publik}
\begin{itemize}
    \item \textbf{Home} -- Landing page dengan informasi masjid
    \item \textbf{Berita} -- Daftar artikel dan detail
    \item \textbf{Galeri} -- Album foto kegiatan
    \item \textbf{Jadwal Kegiatan} -- Kalender interaktif
    \item \textbf{Donasi} -- Form donasi dengan kalkulator zakat
    \item \textbf{Laporan Keuangan} -- Grafik dan tabel keuangan
    \item \textbf{Konsultasi} -- Form tanya jawab
    \item \textbf{Tentang Kami} -- Informasi masjid
\end{itemize}

\subsubsection{Halaman Admin}
\begin{itemize}
    \item \textbf{Dashboard} -- Ringkasan statistik
    \item \textbf{Manajemen User} -- CRUD pengguna
    \item \textbf{Manajemen Role} -- CRUD role \& permission
    \item \textbf{Postingan} -- CRUD artikel + approval
    \item \textbf{Galeri} -- CRUD album foto
    \item \textbf{Kegiatan} -- CRUD jadwal
    \item \textbf{Keuangan} -- Input transaksi + grafik
    \item \textbf{Donasi} -- Verifikasi konfirmasi
    \item \textbf{Konsultasi} -- Chat dengan jamaah
    \item \textbf{Lost \& Found} -- Manajemen barang
    \item \textbf{Form Builder} -- Buat form dinamis
    \item \textbf{Website Info} -- Edit halaman statis
\end{itemize}

\subsection{Screenshot Aplikasi}

\textit{[Tambahkan screenshot aplikasi di sini]}

\vspace{5cm}

\textit{Screenshot 1: Halaman Home}

\vspace{5cm}

\textit{Screenshot 2: Dashboard Admin}

\vspace{5cm}

\textit{Screenshot 3: Halaman Donasi}

\newpage

% ============ BAB 5 ============
\section{RANCANGAN MODUL (Detailed Design)}

\subsection{Modul Autentikasi}

\subsubsection{Spesifikasi: Login}
\begin{table}[h]
\centering
\begin{tabular}{|l|l|}
\hline
\textbf{Aspek} & \textbf{Detail} \\
\hline
Controller & AuthController \\
Method & login(Request \$request) \\
Route & POST /login \\
\hline
\end{tabular}
\end{table}

\textbf{Input:}
\begin{table}[h]
\centering
\begin{tabular}{|l|l|l|}
\hline
\textbf{Field} & \textbf{Tipe} & \textbf{Validasi} \\
\hline
email & string & required, email \\
password & string & required \\
remember & boolean & optional \\
\hline
\end{tabular}
\end{table}

\textbf{Output:}
\begin{table}[h]
\centering
\begin{tabular}{|l|l|}
\hline
\textbf{Kondisi} & \textbf{Response} \\
\hline
Sukses (terverifikasi) & Redirect ke /admin atau / \\
Sukses (belum verifikasi) & Redirect ke /auth/verify \\
Gagal & Kembali dengan error \\
\hline
\end{tabular}
\end{table}

\textbf{Deskripsi:}
Validasi kredensial pengguna, cek verifikasi email, buat session jika berhasil.

\subsubsection{Spesifikasi: Register}

\textbf{Input:}
\begin{table}[h]
\centering
\begin{tabular}{|l|l|l|}
\hline
\textbf{Field} & \textbf{Tipe} & \textbf{Validasi} \\
\hline
full\_name & string & required, max:255 \\
username & string & required, unique \\
email & string & required, email, unique \\
phone\_number & string & required \\
password & string & required, min:8, confirmed \\
captcha & string & required, valid captcha \\
\hline
\end{tabular}
\end{table}

\subsection{Modul Postingan}

\subsubsection{Spesifikasi: Tambah Postingan}

\textbf{Input:}
\begin{table}[h]
\centering
\begin{tabular}{|l|l|l|}
\hline
\textbf{Field} & \textbf{Tipe} & \textbf{Validasi} \\
\hline
judul & string & required, max:255 \\
featured\_image & file & required, image, max:5MB \\
kontent & string & required (HTML dari Quill) \\
\hline
\end{tabular}
\end{table}

\textbf{Proses:}
\begin{enumerate}
    \item Validasi input
    \item Generate slug dari judul
    \item Upload featured image ke storage
    \item Process gambar base64 dalam konten Quill
    \item Simpan ke database dengan status pending
    \item Redirect ke daftar postingan
\end{enumerate}

\subsection{Modul Donasi}

\subsubsection{Spesifikasi: Hitung Donasi}

\textbf{Input:}
\begin{table}[h]
\centering
\begin{tabular}{|l|l|l|}
\hline
\textbf{Field} & \textbf{Tipe} & \textbf{Validasi} \\
\hline
donation\_category & string & required, in:zakat,infaq \\
donation\_type & string & required \\
bank\_id & integer & required, exists:bank\_accounts \\
infaq\_amount & numeric & required if infaq, min:10000 \\
\hline
\end{tabular}
\end{table}

\textbf{Proses:}
\begin{enumerate}
    \item Validasi kategori donasi (zakat/infaq)
    \item Jika zakat: hitung menggunakan ZakatService
    \item Jika infaq: ambil nominal langsung
    \item Simpan hasil ke session
    \item Redirect ke halaman konfirmasi
\end{enumerate}

\subsection{Modul Konsultasi}

\subsubsection{Spesifikasi: Buat Konsultasi}

\textbf{Input:}
\begin{table}[h]
\centering
\begin{tabular}{|l|l|l|}
\hline
\textbf{Field} & \textbf{Tipe} & \textbf{Validasi} \\
\hline
question\_subject & string & required, max:255 \\
question\_text & string & required, min:10 \\
is\_anonymous & boolean & optional \\
\hline
\end{tabular}
\end{table}

\textbf{Output (JSON):}
\begin{table}[h]
\centering
\begin{tabular}{|l|l|}
\hline
\textbf{Kondisi} & \textbf{Response} \\
\hline
Sukses & \{success: true, redirect: url\} \\
Sudah ada pending & Error 422 \\
Belum verifikasi & Error 403 \\
\hline
\end{tabular}
\end{table}

\subsection{Modul Keuangan}

\subsubsection{Spesifikasi: Tambah Transaksi}

\textbf{Input:}
\begin{table}[h]
\centering
\begin{tabular}{|l|l|l|}
\hline
\textbf{Field} & \textbf{Tipe} & \textbf{Validasi} \\
\hline
type & string & required, in:pemasukan,pengeluaran \\
bank\_name & string & required \\
category & string & required \\
transaction\_date & date & required \\
amount & numeric & required, min:1000 \\
proof\_file & file & required, mimes:jpeg,png,jpg,pdf \\
description & string & optional \\
\hline
\end{tabular}
\end{table}

\newpage

% ============ BAB 6 ============
\section{MODEL IMPLEMENTASI}

\subsection{Implementasi Model User}
\begin{lstlisting}[language=PHP]
// app/Models/User.php
class User extends Authenticatable implements MustVerifyEmail
{
    protected $fillable = [
        'role_id', 'username', 'full_name', 'phone_number',
        'image_url', 'email', 'email_verified_at', 'password',
    ];

    public function role()
    {
        return $this->belongsTo(Role::class);
    }

    public function hasPermission($permission): bool
    {
        return $this->role?->permissions()
            ->where('name', $permission)->exists() ?? false;
    }

    public function hasRole($roleName): bool
    {
        return $this->role && 
            strtolower($this->role->name) == strtolower($roleName);
    }
}
\end{lstlisting}

\subsection{Implementasi Middleware Permission}
\begin{lstlisting}[language=PHP]
// app/Http/Middleware/CheckPermission.php
class CheckPermission
{
    public function handle(Request $request, Closure $next, 
        $permissions): Response
    {
        if (!Auth::check()) {
            return redirect('login');
        }
        
        $user = Auth::user();
        $permissions = explode('|', $permissions);
        
        foreach ($permissions as $permission) {
            if ($user->hasPermission($permission)) {
                return $next($request);
            }
        }

        abort(403, 'Unauthorized action.');
    }
}
\end{lstlisting}

\subsection{Implementasi Model Role \& Permission}
\begin{lstlisting}[language=PHP]
// app/Models/Role.php
class Role extends Model
{
    protected $fillable = ['name', 'alias', 'description'];

    public function permissions()
    {
        return $this->belongsToMany(Permission::class, 
            'role_permission');
    }

    public function users()
    {
        return $this->hasMany(User::class);
    }
}

// app/Models/Permission.php
class Permission extends Model
{
    protected $fillable = ['name', 'group'];

    public function roles()
    {
        return $this->belongsToMany(Role::class, 
            'role_permission');
    }
}
\end{lstlisting}

\newpage

% ============ BAB 7 ============
\section{PENGUJIAN}

\subsection{Skenario Pengujian Modul Autentikasi}
\begin{longtable}{|c|l|p{3cm}|p{3cm}|p{3cm}|}
\hline
\textbf{No} & \textbf{Fitur} & \textbf{Data Test} & \textbf{Skenario} & \textbf{Expected} \\
\hline
\endfirsthead
1 & Login & Email valid, Pass valid & Login valid & Redirect dashboard \\
2 & Login & Email valid, Pass salah & Login gagal & Error message \\
3 & Login & Email tidak ada & Login gagal & Error message \\
4 & Register & Data lengkap & Registrasi & Redirect verify \\
5 & Register & Email duplikat & Registrasi gagal & Error email \\
6 & OTP & Kode benar & Verifikasi & Email terverifikasi \\
7 & OTP & Kode salah & Verifikasi gagal & Error kode \\
\hline
\caption{Skenario Pengujian Autentikasi}
\end{longtable}

\subsection{Skenario Pengujian Modul Postingan}
\begin{longtable}{|c|l|p{3cm}|p{3cm}|p{3cm}|}
\hline
\textbf{No} & \textbf{Fitur} & \textbf{Data Test} & \textbf{Skenario} & \textbf{Expected} \\
\hline
1 & Create & Data valid & Tambah postingan & Status pending \\
2 & Create & Tanpa gambar & Tambah gagal & Error validasi \\
3 & Approval & action=approve & Approve & Status published \\
4 & Approval & action=reject & Reject & Status rejected \\
5 & Delete & ID valid & Hapus & Data terhapus \\
\hline
\caption{Skenario Pengujian Postingan}
\end{longtable}

\subsection{Skenario Pengujian Modul Donasi}
\begin{longtable}{|c|l|p{3cm}|p{3cm}|p{3cm}|}
\hline
\textbf{No} & \textbf{Fitur} & \textbf{Data Test} & \textbf{Skenario} & \textbf{Expected} \\
\hline
1 & Zakat Maal & Harta 100jt & Hitung & Hasil 2.25jt \\
2 & Zakat Fitrah & 3 jiwa & Hitung & Hasil 135rb \\
3 & Infaq & 50rb & Donasi & Tersimpan \\
4 & Konfirmasi & Bukti valid & Konfirmasi & Status Pending \\
5 & Verify & ID valid & Approve & Status Verified \\
\hline
\caption{Skenario Pengujian Donasi}
\end{longtable}

\subsection{Skenario Pengujian Modul Keuangan}
\begin{longtable}{|c|l|p{3cm}|p{3cm}|p{3cm}|}
\hline
\textbf{No} & \textbf{Fitur} & \textbf{Data Test} & \textbf{Skenario} & \textbf{Expected} \\
\hline
1 & Pemasukan & 1jt, bukti valid & Input & Saldo bertambah \\
2 & Pengeluaran & 500rb, saldo cukup & Input & Tersimpan \\
3 & Pengeluaran & Melebihi saldo & Input gagal & Error saldo \\
4 & Hapus & ID valid & Hapus & Terhapus \\
\hline
\caption{Skenario Pengujian Keuangan}
\end{longtable}

\subsection{Skenario Pengujian Akses Permission}
\begin{longtable}{|c|l|l|l|c|}
\hline
\textbf{No} & \textbf{Role} & \textbf{Route} & \textbf{Permission} & \textbf{Result} \\
\hline
1 & Humas & /admin/postingan & view\_posts & Allowed \\
2 & Humas & /admin/keuangan & view\_finance & 403 \\
3 & Bendahara & /admin/keuangan & view\_finance & Allowed \\
4 & Bendahara & /admin/postingan & view\_posts & 403 \\
5 & Sarpras & /admin/barang-hilang & view\_lost\_items & Allowed \\
6 & Jamaah & /admin/* & view\_dashboard & 403 \\
\hline
\caption{Skenario Pengujian Permission}
\end{longtable}

\newpage

% ============ LAMPIRAN ============
\appendix
\section{Daftar Permission}
\begin{longtable}{|l|l|p{5cm}|}
\hline
\textbf{Permission} & \textbf{Group} & \textbf{Deskripsi} \\
\hline
view\_dashboard & Dashboard & Akses dashboard admin \\
view\_roles & Role & Lihat daftar role \\
create\_roles & Role & Buat role baru \\
edit\_roles & Role & Edit role \\
delete\_roles & Role & Hapus role \\
view\_users & Pengguna & Lihat daftar pengguna \\
edit\_users & Pengguna & Edit pengguna \\
delete\_users & Pengguna & Hapus pengguna \\
view\_consultations & Konsultasi & Lihat konsultasi \\
reply\_consultations & Konsultasi & Balas konsultasi \\
view\_events & Kegiatan & Lihat kegiatan \\
create\_events & Kegiatan & Buat kegiatan \\
view\_posts & Postingan & Lihat postingan \\
create\_posts & Postingan & Buat postingan \\
approve\_posts & Postingan & Setujui postingan \\
view\_gallery & Galeri & Lihat galeri \\
view\_finance & Keuangan & Lihat keuangan \\
manage\_income & Keuangan & Kelola pemasukan \\
verify\_donation & Keuangan & Verifikasi donasi \\
view\_lost\_items & Barang Hilang & Lihat barang hilang \\
\hline
\caption{Daftar Permission Sistem}
\end{longtable}

\section{Akun Demo}
\begin{table}[h]
\centering
\begin{tabular}{|l|l|l|l|}
\hline
\textbf{Role} & \textbf{Username} & \textbf{Email} & \textbf{Password} \\
\hline
Super Admin & superadmin & superadmin@samak.com & password123 \\
Admin & admin & admin@samak.com & password123 \\
Humas & humas & humas@samak.com & password123 \\
Bendahara M & bendaharaPemasukan & bendaharaM@samak.com & password123 \\
Bendahara K & bendaharaPengeluaran & bendaharaK@samak.com & password123 \\
Sarpras & sarpras & sarpras@samak.com & password123 \\
Jamaah & jamaah & jamaah@samak.com & password123 \\
Koor Humas & koordinator\_humas & koor\_humas@samak.com & password123 \\
\hline
\end{tabular}
\caption{Daftar Akun Demo}
\end{table}

\vfill
\begin{center}
\textit{Dokumen ini dibuat berdasarkan analisis kode sumber proyek Samak Pro}
\end{center}

\end{document}
